\section{Conclusions}
\label{sec:conclusions}

This paper asked whether any of the nearby stars that ALMA has observed shows
a narrowband, drifting radio carrier of the kind an artificial transmitter
would produce at millimetre and submillimetre wavelengths. Because the
observations were taken for other purposes, it also asked how deep and how
complete a search of the archive for such a carrier can be. To answer both
we searched the archival ALMA spectra of \NSysClassA{} star systems within
40\,pc. The fine-channel windows cover \CvUnionA\,GHz of sky frequency, in
\CovNIslands{} separate intervals between \SurvFreqLoA{} and
\SurvFreqHiA\,GHz. That union is \SbPriorRatio{} times the \SbPriorBandMHz\,MHz of the
wider of the two millimetre searches of nearby stars that preceded it, and
the coverage of the median system alone is still \SbPriorRatioSysMed{}
times larger. We removed the drift from each spectrum and compared it with
control positions in the same field, so that the background at the star
could be judged from the observation itself.

No transmitter was found. In \CsNCross{} windows the stellar statistic reaches
the threshold $T_\star\geq\StageNullTrig$. Of these, \MkNCoinc{} coincide with
catalogued molecular lines, \LnAttrSupp{} of them with further evidence of
line emission, and none was identified as interference by the available
diagnostics. Counted by the adopted rule, each window by its largest cell, the
other \CeObs{} are consistent with chance. The expected number, built
window by window from the control positions beside each star, is \CeExp, with
$P=\CeProbNull$. Built instead from the control positions of all the windows
pooled together, as a check, it is \CeExpPool, with $P=\CeProbNullPool$. The extraction field and the line
mask could each have been chosen in two ways. Corrected for having been
picked out of those \CgNTrialWord{} combinations, the smallest of their
probabilities is \CgBonfPhrase, \CgVerdictClause. That correction allows only
for the choice among those combinations, and it does not make good every
decision taken after the data had been seen. All of these probabilities are
conditional on the line mask as adopted, which was completed
after the crossings had been examined, and on the control positions being a
fair model of the background at the star. They also depend on the counting
rule, which is applied at the star and at the control positions alike.
Counting every unattributed cell instead gives
\AcObsCell{} against \AcExpCell{} expected, with $P=\AcProbCell$, or
\AcProbCellTwoPhrase{} when doubled for the choice between the two ways of
counting cells and \AcProbCellThree{} when trebled for the three rules
compared. The excess lies in \AcStarSecNWord{} weaker cells behind the largest ones,
where the control positions predict about \AcSecExp, and none of them is
found again in another block.

None of the unattributed crossings became a recurrent feature, although
the molecular lines did, as they should. Wherever an unattributed
crossing was observed again in a separated session and the repeat could be
read, nothing was found at its own frequency in the star's frame. That shows
only that none of them returns in the way a line would. A drifting carrier
would by then have moved, and the archive can test one only by asking
whether a single trajectory joins two crossings, which can tell crossings
apart but can never confirm them. Recurrence is in any case necessary for an
artificial origin but not sufficient: a feature found again would still need
its line identity, interference and the behaviour of its frequency in the
star's frame to be checked. For the crossings themselves, therefore,
this is a non-detection, not an exclusion.

Because nothing was confirmed, what the search delivers is a set of
single-epoch upper limits. In the \XwNWin{} windows where an unattributed
crossing, or a weaker unattributed cell, is already present, a carrier
reaching only the threshold would not have stood out, so there the limit is
quoted at the statistic actually observed, which errs on the conservative
side. That moves the limits of \XwNSysWord{} systems, marked in
Table~\ref{tab:persys}, and \XwHeadlineClause. The limits apply to a
carrier narrower than one channel, lying outside the molecular-line mask,
radiating while ALMA observed and drifting at a constant rate of no more
than \DriftCeilLo--\DriftCeilHi\,nHz, the drift ceiling discussed in
\S\ref{sec:excludes}. The injected carriers that calibrate them drift
linearly, so a track that curves during the integration is not covered.
To reach the threshold nine times in ten in a single epoch, such a carrier needs an equivalent isotropic radiated power, or EIRP,
with a median, over the \NSysALimSum{} systems with a bounded limit, of
$\EirpNinetySysMedA$\,W in each system's most sensitive window. These limits
are corrected for the atmospheric decorrelation wherever the phase rms was
measured in the block, and elsewhere by a value imputed from blocks matched in
baseline length, which are marked in Table~\ref{tab:persys}. The phase rms is
measured on the calibrator and not along the line of sight to the star, so
the correction is itself a floor, and the limits remain optimistic by any
loss it misses. This is a single-epoch trigger completeness, not the
chance of confirming a signal, since a crossing would still have to be found
again. The limit is not the same at every frequency a system was observed
at, and Fig.~\ref{fig:waterfall}c shows how much of the covered band each
power reaches. At $\CovEirpRefTex$\,W it is \CovUnionAtRef\,GHz of the
\CvUnionA\,GHz searched. The deepest limit, $\EirpNinetySysLoA$\,W, is reached toward
$\epsilon$~Eridani, and \NSysEfifteen{} of the \NSysALimSum{} systems with a
limit that does not depend on an unchecked position reach
$\CovEirpRefTex$\,W. $\epsilon$~Eridani, however, lies
\EdDeepOffLo--\EdDeepOffHi$''$ from the phase centre, where the primary-beam
response is below \OffAcut, and no carrier was injected further out than
\OffInjMaxArcsec$''$, so its limit rests on the model of the beam alone. The deepest limit that does rest on injected
configurations is $\PtDeepInjEirp$\,W, toward \PtDeepInjStar. Nor is
$\epsilon$~Eridani's limit corrected for decorrelation, since neither its block
nor any block matched to it records the phase rms. Its spectrum was extracted at
the star's apparent position, so the parallax is allowed for.

It helps to set these numbers against something we have built ourselves.
The Arecibo planetary radar, one of the most powerful transmitters ever
operated on the Earth, had an EIRP of about $\AroEirp$\,W at
\AroFreqGHz\,GHz. The median limit found here is about \AroMedMult{} times
that, and even the deepest is \AroDeepMult{} times it, so
\AroNSysBelowWord{} system reaches an Arecibo-like transmitter, although
\AroNSysTen{} come within a factor of ten of it (counted on the unrounded
limits, \S\ref{sec:mmdifferent}). The same comparison can be
put as a distance. Because the received flux falls as the square of the
distance, an Arecibo-like transmitter placed at a star of this survey would
be detected only if the system were much closer to us than it is. For the
median system it would have to be within \AroDistMedPc\,pc, and even in the
most favourable case, \AroDistMaxStar, whose limit is deep for its distance,
within \AroDistMaxPc\,pc, a value whose decorrelation correction is
\AroDistMaxBasis. A \BenchDiam\,m dish radiating \BenchPowerMW\,MW is
a more natural transmitter at these frequencies. Tuned to the lower edge of
each system's most sensitive window, between \BenchTunedFreqLo{} and
\BenchTunedFreqHi\,GHz, its EIRP lies above the single-epoch limit of
\BenchNSysTuned{} of the \NSysALimSum{} systems. It could only be seen, of
course, if its beam happened to point at us, as discussed in
\S\ref{sec:mmdifferent}.

These limits come from the single-epoch carrier search, and they are not the
deepest number in the paper. Toward Proxima Centauri the Class~B spectra,
stacked in the star's own frame, reach $\SzProxEirpCensus$\,W in
\ChanBModeMHz\,MHz channels (\S\ref{sec:stack}). That is a limit of a
different kind. It applies to an unresolved spectral excess that holds its
frequency in the stellar frame to within one of those wide channels over all
the epochs combined, it carries its own injection calibration, and it cannot
tell a carrier from a line. A transmitter on \PxPlanet{} would move across
many channels in its orbit, and unless the orbit is seen nearly face-on the
stack would dilute it by roughly the number of channels crossed, about
\PxNChanWalk{} at $\sin i=1$. For a drifting carrier the single-epoch
limits therefore remain the ones that apply.

Confirmation is limited by repeat coverage, by the depth of
the repeats that do exist and, for a drifting carrier, by a repeat
frequency that is uncertain by a median of \ZtRegChanMed{} channels. A
crossing can be put to the test only if its star is observed again in a
separated session, at the frequency the crossing predicts. Of the
\NSysClassA{} systems, \ZcNoRepNoSep{} were never observed in two separated
sessions at all, so that nothing found toward them could have been confirmed
at any power. For \ZcNoRepNoFreqWord{} more, separated sessions exist but
share no covered frequency, so the same-frequency recurrence test cannot be
applied there. Where a repeat exists it may still be too shallow, as it was
for \HjNPowerLo{} of the \RcNExcl{} crossings read at their predicted cell.
A drifting carrier, finally, would have to be looked for again at a
frequency carried forward from its own fitted drift, and one block does not
measure the drift well enough to say where that will be. A programme designed for the purpose, on the same telescope,
would use fine channels throughout, tune away from the bright molecular
lines and schedule its repeat observations deliberately, so that any
crossing it found could be looked for again at the frequency it predicts
(\S\ref{sec:future}).

What, then, can be said about transmitters? Take a continuous carrier
narrower than a channel at one of these stars, outside the line mask and
drifting at a constant rate within the ceiling. At or above the system's limit, and within the
frequency range of its most sensitive window, which lies between
\NtBestNuLo{} and \NtBestNuHi\,GHz across the systems, it would have
produced a threshold crossing in nine trials in ten. That is a completeness and not a
deterministic exclusion. Of \NSysALimSum{} carriers, one sitting exactly at
each system's limit, about \NtMissAtLimWord{} would have been missed, and
more if the correction for decorrelation falls short.

A great deal remains possible. A transmitter below these limits would not have been seen,
and neither would one that was switched off, or whose beam was not pointed
at the Earth, while ALMA was looking. Signals at frequencies between the
covered intervals or inside the line mask, or drifting faster than the
ceiling, are not constrained at all, and signals broader than a channel only
weakly. Nor is anything said
about the nearby stars that ALMA has not observed, which are almost all of
those within 40\,pc.

%% ====================================================================
%% HOW THIS SECTION IS BUILT (2026-11-09).  It runs as the abstract does:
%% the question (para 1), what was found (paras 2-3: crossings, chance as
%% conditional on the adopted mask and the control-position background,
%% expected count built window by window; the trials factor for the four
%% field/mask combinations, which covers that one choice and no other
%% post-hoc decision; no recurrence; a non-detection, not an exclusion),
%% the single-epoch limits (paras 4-6: ninety per cent trigger completeness
%% for a CONSTANT-DRIFT carrier (the injections drift linearly), the median
%% over the systems with a bounded limit, the limits CORRECTED for
%% decorrelation where the phase rms was measured, that correction a floor
%% so the limits stay optimistic by any residual; how the limit varies with
%% frequency, Fig. 2c; the deepest value, epsilon Eri, beyond the injected
%% offsets, with the deepest limit resting on injected configurations
%% beside it; Arecibo and the 12 m benchmark on the corrected limits; the
%% Class B stack, which DILUTES a drifting carrier by the channels
%% crossed), what limits confirmation (para 7: ONE sentence names all
%% three -- repeat coverage, the depth of the repeats, and for a drifting
%% carrier a repeat frequency uncertain by ZtRegChanMed channels -- then
%% each in turn; the old "coverage, not depth" is withdrawn),
%% and what the survey would have caught and what remains possible (paras
%% 8-9).  NEVER "excluded" for transmitters: the survey would have produced
%% a crossing in nine trials in ten, which is a completeness, not a
%% deterministic exclusion.  No citations.  Macro names are written here
%% without backslashes on purpose.
%% ====================================================================
